\chapter{Licences}\label{cha:licences}\markboth{Licences}{Licences}

The Thank You chapter (page~\pageref{cha:thanks}) is the thanks; this
is the record. The
authoritative copy is \pth{LICENSES.md} in the repository --- if the
two ever disagree, that file wins.

\section{The project itself}
The K4510 as a whole is distributed under the \textbf{GNU General
Public License, version~2 or (at your option) any later version}.
Everything written for this project --- the machine, VICKY, SHEILA,
the DMA engine, the MATH unit, the system ROM and K/OS, JIM, the Tube
ULA, the network layer, the editors, the demos, RX and the Pascal
target --- is Copyright \copyright{} 2026 Michael
Longval and is released under those terms. The choice is not really a
choice: the CPU core and the OPL2 emulation are GPL-2.0-or-later, and a
work containing them must be too.

\section{Code inside the machine}
\begin{itemize}[leftmargin=1.4em]
\item \textbf{Xemu 65xx/45GS02 CPU core} --- Gábor Lénárt.
  \pth{core/xemu/}: \verb|cpu65.c|, \verb|cpu65.h|,
  \pth{cpu65_mega65_timings.h}, \pth{cpu65ce02_disasm_tables.c},
  used unchanged. \emph{GPL-2.0-or-later.}
\item \textbf{fmopl (YM3812/OPL2)} --- Jarek Burczyński and Tatsuyuki
  Satoh, from MAME by way of VICE. \verb|core/opl2/|, vendored
  unaltered. \emph{GPL-2.0-or-later.}
\item \textbf{reSID} --- Dag Lem, as shipped in VICE~3.3, was
  \verb|core/resid/| until 2026-09-05 and is no longer in the tree
  (Appendix~\ref{cha:sound}). \emph{GPL-2.0-or-later}; recorded here
  because a vendored component that quietly disappears is worse than one
  recorded as gone.
\item \textbf{BBC BASIC for SDL 2.0, console edition (BBCTTY)} ---
  Richard T. Russell. \verb|tube/|, vendored and altered (see
  \pth{tube/ALTERED.md}). \emph{zlib licence.} ``BBC BASIC'' is the
  name of Mr Russell's interpreter, which he publishes under his own
  arrangement with the BBC. This project has not sought, and does not
  hold, any licence to the name: it is used here only to identify the
  interpreter, which is why this book says ``Richard Russell's BBC
  BASIC'' and not ``the K4510's''.
\item \textbf{EhBASIC 2.22} --- Lee Davison, in the ca65 form from
  jefftranter/6502. \verb|basic/basic.asm|. \emph{Free for
  non-commercial use only}; derivatives must carry the words
  ``Derived from EhBASIC'' --- see \pth{basic/README-EhBASIC.txt}.
  It ships as a separate program
  (\pth{fs/LANG/EHBASIC/ehbasic.prg}) and is not linked with the GPL
  code.
\item \textbf{RunCPM} --- Marcelo Dantas. \verb|cpm/src/|, vendored
  unmodified (\pth{cpm/VENDORED-FROM.txt}). \emph{MIT.}
\item \textbf{Tali Forth 2} --- Scot W. Stevenson, Sam Colwell,
  Patrick Surry. \pth{forth/tali/}, vendored unmodified.
  \emph{Public domain.}
\item \textbf{Wozmon} --- Steve Wozniak's 1976 monitor, reimplemented
  here for the 45GS10. Apple published the original listing; this
  machine's version is the project's own code under the project
  licence.
\item \textbf{cc65 runtime} --- the ROM and every \verb|.prg| are
  linked against \verb|none.lib|. \emph{cc65's zlib-style licence.}
\item \textbf{Microsoft 6502 BASIC} --- \emph{MIT}, released by
  Microsoft in 2025, via \href{https://github.com/mist64/msbasic}{mist64/msbasic}
  at a pure-Microsoft configuration. \pth{basic/msbasic/}, vendored
  unmodified --- and only the files a pure build assembles: none of the
  per-manufacturer material is here, which is the line that keeps the
  build resting on the MIT release alone.
\item \textbf{Supermon+64 1.2} --- Jim Butterfield, restored and
  commented by J.\,B. Langston
  (\href{https://github.com/jblang/supermon64}{jblang/supermon64}),
  ported to the 45GS10. Butterfield published it for anyone to use;
  Langston asks for attribution.
\end{itemize}

\section{Fonts}
\begin{itemize}[leftmargin=1.4em]
\item \textbf{Linux kernel 8x8 console font} --- \pth{data/font8.bin},
  built by \pth{data/mkfont.py} from \pth{lib/fonts/font_8x8.c}.
  \emph{GPL-2.0.} This is the default text font.
\item \textbf{MEGA65 open-roms chargen + PXLfont 2.3} ---
  \pth{data/fonts/openroms/}. Clean-room C64/C65 character shapes;
  PXLfont is Retrofan's, included with open-roms' recorded permission.
  \emph{LGPL-3.0-or-later} --- \pth{8x8font.png} is shipped as the
  editable source, for LGPL compliance.
\item \textbf{unscii-8} --- Ville-Matias Heikkilä.
  \pth{data/fonts/unscii/}; \pth{font8-unscii.bin} is a generated
  drop-in alternative to \pth{data/font8.bin}. \emph{Public domain.}
\item \textbf{BESCII v3} (Mono + source) --- Damian Vila.
  \pth{data/fonts/bescii/}. \emph{CC0-1.0.}
\item \textbf{Clear Sans} (Intel, \emph{Apache-2.0}) and
  \textbf{Iosevka} (Belleve Invis, \emph{SIL OFL 1.1}) --- the two
  faces this book is set in. They are not part of the machine.
\end{itemize}

\begin{aside}
\textbf{No Commodore ROMs here.} Until 2026-08-23 the text font was
derived from a Commodore~64 character ROM. That file and every
derivative were removed from the repository \emph{and from its
history} before it was made public. If you own a C64 and want the real
thing, drop your own \verb|chargen.bin| into \pth{/SYSTEM/ETC}: the
machine will use it, and \verb|.gitignore| will keep it out of the
repository. That copy is yours, not ours to publish.
\end{aside}

\section{The Linux the machine boots on}
The live image is a Debian, assembled by \pth{linux/build-live.sh} from
Debian's own archive, so every package in it carries its own terms and
Debian's copyright files travel with it. One
thing is built from source into it:
\begin{itemize}[leftmargin=1.4em]
\item \textbf{Tek40xx} --- Ian Schofield's Tektronix 4010/4014
  terminal, built from upstream with one patch of ours.
  \emph{GPL-3.0.} It is a separate program on the Linux side, not part
  of the emulator.
\end{itemize}
\noindent The Tektronix plot files in \pth{linux/tek40xx/plt/} are
Tek40xx's own gnuplot examples plus a page generated here. The vintage
cassette plots that circulate with \emph{other} Tektronix emulators are
marked personal-use-only by their authors and are deliberately \emph{not}
vendored; the README says how to fetch them for yourself.

\begin{aside}
\textbf{The Raspberry Pi build} was here until 2026-09-07 and needed two
components that never lived in this repository: \textbf{Circle} (Rene
Stange, \emph{GPL-3.0}) and \textbf{circle-libsdl2} (Xalior,
\emph{zlib}, with a \emph{GPL-3.0} linker script). A \verb|kernel8.img|
built from those was GPL-3.0 as a whole, which GPL-2.0-\emph{or-later}
code permits. The tag \texttt{alpha-0.5} is the last release that
shipped one.
\end{aside}

\section{Build tools}
Not distributed with the machine, but required to build it: GCC and
GNU Make (\emph{GPL}), cc65 (\emph{zlib-style}), 64tass
(\emph{GPL-2.0}), NASM (\emph{BSD-2-Clause}), SDL2 (\emph{zlib}),
Python~3 (\emph{PSF}), and for this book XeLaTeX (\emph{LPPL}). Mad
Pascal and MADS (Tomasz Biela, \emph{MIT}) are needed only if you
compile Pascal; the K4510 target in \verb|pascal/| is installed into
your own Mad Pascal checkout.

\section{What is deliberately not shipped}
\begin{itemize}[leftmargin=1.4em]
\item \textbf{The CP/M system disk.} RunCPM's \verb|DISK/A0.zip| ---
  Digital Research's \verb|ASM|, \verb|MAC|, \verb|DDT|, \verb|STAT|,
  \verb|SUBMIT| and third-party tools --- installs into
  \pth{fs/CPM/A/0/} for your use but is not committed: its files have
  mixed provenance and this repository is public. The same goes for
  WordStar, Turbo Pascal~3 and MBASIC, which the \verb|K-*.SUB|
  launchers know how to start but which you must supply yourself.
\item \textbf{ZX Origins fonts.} Damien Guard's faces are free to use
  and \emph{not} free to re-host, so the machine names them in the F7
  font menu and does not ship them: \pth{tools/mkzxfonts.py} turns your
  own download into screen fonts, and without it those entries fall back
  to the default face.
\item \textbf{A Commodore character ROM} --- see the note above.
\item \textbf{Your \texttt{STARTUP.BAT}} --- yours, not the repository's
  (copy \pth{/SYSTEM/ETC/STARTUP.SAMPLE} to start one).
\end{itemize}

\section{If you redistribute the machine}
Ship the sources you built from, keep \pth{LICENSES.md} and this
chapter with them, keep \pth{8x8font.png} beside the open-roms font,
and remember two things that are easy to forget: EhBASIC is
non-commercial-only, and ``BBC BASIC'' is the name of Richard Russell's
interpreter rather than of anything made here --- describe it as his,
as this book does, and give whatever you build on it a name of its own.
